\documentclass[11pt,a4paper]{article}
\usepackage[utf8]{inputenc}
\usepackage[T1]{fontenc}
\usepackage[french]{babel}
\frenchbsetup{StandardLists=true}
\usepackage[left=2cm,right=2cm,top=2cm,bottom=2cm]{geometry}
\usepackage[dvipsnames]{xcolor}
\usepackage{fouriernc}
\usepackage{amsmath,amssymb,amsfonts}
\usepackage{array,tabularx,booktabs,multirow}
\usepackage{colortbl}
\usepackage{graphicx}
\usepackage{mdframed}
\usepackage{enumitem}
\usepackage{fancyhdr}
\usepackage{chemformula}
\usepackage{awesomebox}
\setlength{\aweboxvskip}{0mm}
\usepackage{tcolorbox}
\tcbuselibrary{skins}
\usepackage{fontawesome5}

\definecolor{exoheader}{HTML}{1E5AA0}
\definecolor{exolight}{HTML}{DCE8FF}
\definecolor{warnorange}{RGB}{220,100,0}
\definecolor{problemebg}{HTML}{FFFDE7}

\renewcommand{\arraystretch}{1.8}
\setlength{\extrarowheight}{1mm}

\pagestyle{fancy}
\renewcommand\headrulewidth{1pt}
\fancyhead[L]{Académie de Besançon}
\fancyhead[C]{\textbf{TP --- Dosage conductimétrique (ExAO)}}
\fancyhead[R]{Terminale / BTS Chimie}
\fancyfoot[C]{page \thepage}
\fancyfoot[R]{2025 -- 2026}
\fancyfoot[L]{Alexis RUIZ}

\newcommand{\sectitle}[1]{%
  \vspace{6pt}
  \noindent\colorbox{exoheader}{\parbox{\dimexpr\linewidth-2\fboxsep}%
    {\color{white}\bfseries\large\strut\quad #1}}
  \vspace{4pt}
}
\newcommand{\subsectitle}[1]{%
  \vspace{4pt}
  \noindent\colorbox{exolight}{\parbox{\dimexpr\linewidth-2\fboxsep}%
    {\color{exoheader}\bfseries\strut\quad #1}}
  \vspace{2pt}
}
\newcounter{question}
\newcommand{\question}{%
  \stepcounter{question}%
  \noindent\textcolor{exoheader}{\textbf{Question \thequestion.}}\quad
}

\begin{document}

\noindent\textbf{NOM :}\hfill\textbf{Prénom :}\hfill\textbf{Classe :}\hfill\phantom{x}\\[0.2cm]
\hrule\vspace{0.2cm}

\begin{center}
  {\large\bfseries\textcolor{exoheader}{Travaux Pratiques de Chimie}}\\[4pt]
  {\LARGE\bfseries Dosage conductimétrique du sérum physiologique}\\[2pt]
  {\small Académie de Besançon --- Belfort}
\end{center}
\hrule\vspace{6pt}

\begin{tcolorbox}[enhanced, colback=problemebg, colframe=exoheader, boxrule=1.5pt,
  title={\faFlask\quad Problématique},
  fonttitle=\bfseries\color{white},
  attach boxed title to top left={yshift=-2mm, xshift=4mm},
  boxed title style={colback=exoheader, colframe=exoheader}]
\textit{Le sérum physiologique est une solution de chlorure de sodium à 0,9\% en masse.
\textbf{Peut-on vérifier cette concentration par dosage conductimétrique avec courbe d'étalonnage ?}}
\end{tcolorbox}

\vspace{6pt}

% ══════════════════════════════════════════════════════════════════════════════
\sectitle{Partie I --- Rappels théoriques}
% ══════════════════════════════════════════════════════════════════════════════

\begin{mdframed}[backgroundcolor=exolight, linecolor=exoheader, linewidth=1.2pt]
La \textbf{conductivité} $\sigma$ (en mS/cm) d'une solution ionique est proportionnelle
à la concentration en ions : $\sigma = k \times C$.
La \textbf{courbe d'étalonnage} $\sigma = f(C)$ permet de déterminer la concentration
d'une solution inconnue par lecture graphique.
\end{mdframed}

\vspace{4pt}
\question Quels ions sont présents dans une solution de \ch{NaCl} ? Écrire l'équation de dissolution.

\vspace{0.3cm}\noindent\dotfill

\question Expliquer pourquoi la conductivité augmente avec la concentration en \ch{NaCl}. \dotfill

\vspace{0.3cm}\noindent\dotfill

% ══════════════════════════════════════════════════════════════════════════════
\sectitle{Partie II --- Protocole expérimental}
% ══════════════════════════════════════════════════════════════════════════════

\subsectitle{1. Préparation des solutions étalons de NaCl}

\vspace{4pt}
\begin{center}
\begin{tabularx}{\linewidth}{|>{\centering\bfseries\color{white}\columncolor{exoheader}}m{1.5cm}
                              |>{\centering\arraybackslash}X
                              |>{\centering\arraybackslash}X
                              |>{\centering\arraybackslash}X|}
\hline
\rowcolor{exoheader}
\color{white}\bfseries N° & \color{white}\bfseries $C$ (mol/L) & \color{white}\bfseries $\sigma$ mesurée (mS/cm) & \color{white}\bfseries Remarques \\
\hline
1 & $5{,}0 \times 10^{-4}$ & & \\
\hline
2 & $1{,}0 \times 10^{-3}$ & & \\
\hline
3 & $2{,}0 \times 10^{-3}$ & & \\
\hline
4 & $5{,}0 \times 10^{-3}$ & & \\
\hline
5 & $1{,}0 \times 10^{-2}$ & & \\
\hline
6 & $5{,}0 \times 10^{-3}$ (contrôle) & & \\
\hline
\end{tabularx}
\end{center}

\vspace{4pt}
\subsectitle{2. Mesures avec ExAO}

\begin{enumerate}[leftmargin=1.5cm, label=\textcolor{exoheader}{\textbf{\arabic*.}}]
  \item Connecter la sonde conductimétrique à la console ExAO (voie 1).
  \item Lancer le logiciel \og Atelier Scientifique \fg{} → \og Atelier généraliste \fg{}.
  \item Immerger la sonde dans la solution la \textbf{moins concentrée}.
  \item Lancer l'acquisition ; saisir la concentration ; valider le point.
  \item Répéter pour chaque solution étalon (de la moins concentrée à la plus concentrée).
  \item Arrêter l'acquisition et imprimer la courbe d'étalonnage.
\end{enumerate}

\vspace{4pt}
\subsectitle{3. Préparation de la solution de sérum diluée}

\begin{enumerate}[leftmargin=1.5cm, label=\textcolor{exoheader}{\textbf{\arabic*.}}]
  \item Prélever 10~mL de sérum physiologique à la pipette jaugée.
  \item Transvaser dans une fiole jaugée de 100~mL, compléter à la jauge (dilution 1/10).
  \item Verser dans un bécher étiqueté \og sérum dilué \fg{}.
\end{enumerate}

% ══════════════════════════════════════════════════════════════════════════════
\newpage
\sectitle{Partie III --- Exploitation des résultats}
% ══════════════════════════════════════════════════════════════════════════════

\subsectitle{1. Courbe d'étalonnage}

\question Tracer la courbe d'étalonnage $\sigma = f(C)$ sur le graphique ci-dessous
(ou sur papier millimétré). Vérifier la linéarité.

\vspace{5cm}
\begin{center}
\textit{(graphique $\sigma = f(C)$)}
\end{center}
\vspace{0.5cm}

\subsectitle{2. Concentration du sérum dilué}

\begin{enumerate}[leftmargin=1.5cm, label=\textcolor{exoheader}{\textbf{\arabic*.}}]
  \item Relancer l'acquisition et plonger la sonde dans le sérum dilué.
  \item Mesurer la conductivité : $\sigma_{\text{sérum dilué}} = $ \dotfill~mS/cm
  \item Placer ce point sur la courbe et lire la concentration correspondante.
\end{enumerate}

\[
C_{\text{sérum dilué}} = \dotfill \text{ mol/L}
\]

\subsectitle{3. Concentration du sérum commercial}

\question La solution a été diluée 10 fois. En déduire la concentration molaire du sérum commercial.

\[
C_{\text{sérum}} = 10 \times C_{\text{sérum dilué}} = \dotfill \text{ mol/L}
\]

\question Calculer la concentration massique en \ch{NaCl} du sérum.
($M_{\ch{NaCl}} = 58{,}5$~g/mol)

\[
C_m = C_{\text{sérum}} \times M_{\ch{NaCl}} = \dotfill \text{ g/L}
\]

\question Comparer avec la valeur attendue (9~g/L). Calculer l'écart relatif et conclure.

\vspace{0.3cm}\noindent\dotfill

\vspace{0.3cm}\noindent\dotfill

\end{document}
